\section{The exact short-window half-velocity kernel}
\label{sec:terminal-modal-early-kernel}

The remaining early inequality in
\eqref{eq:terminal-modal-short-early-convolution-open} has a simpler kernel
than either $J_k$ or $J_kL$ separately.  This section records that kernel,
its exact cycle fold, and the cancellation retained from each moving-frontier
source.  The reduction is exact, but the resulting finite signed convolution
is still open.

\begin{lemma}[Half-velocity binomial-tail kernel; proved here]
\label{lem:terminal-modal-half-velocity-kernel}
On the infinite line put
\[
 \mathcal T_k=\left(J_k-\frac12J_{k-1}\right)L,
 \qquad k\ge1,
\]
and write $t_k(r)=[z^r]\mathcal T_k$.  Then
\begin{equation}
 t_k(r)=\frac12\Pr\{\operatorname{Bin}(k,1/2)\ge |r|\},
 \qquad r\in\mathbb Z.
 \label{eq:terminal-modal-half-velocity-tail}
\end{equation}
In particular $t_k(r)\ge0$, $t_k(r)=0$ for $|r|>k$, and
$\max_rt_k(r)=1/2$.

Let
\[
 \bar t_k(r)=\sum_{a\in\mathbb Z}t_k(r+2ma).
\]
For every path vector $d=(d_0,\ldots,d_m)$, even extension and restriction
give the exact folded formula
\begin{equation}
 \begin{split}
 (\mathcal T_kd)(j)={}&\bar t_k(j)d_0+\bar t_k(j-m)d_m\\
 &+\sum_{i=1}^{m-1}
 \{\bar t_k(j-i)+\bar t_k(j+i)\}d_i,
 \qquad0\le j\le m.
 \end{split}
 \label{eq:terminal-modal-half-velocity-fold}
\end{equation}
If $qk<3/50$, then $k<24m/25<m$, so every residue in the displayed
formula has at most one supported $2m$-alias.
\end{lemma}

\begin{proof}
Use the Laurent convention
\[
 a=\frac{1+z}{2},\qquad b=\frac{1+z^{-1}}2,
 \qquad a=zb,qquad L=ab.
\]
The fundamental velocity kernel has the exact form
\begin{equation}
 J_k=b^{k-1}\frac{1-z^k}{1-z}.
 \label{eq:terminal-modal-J-Laurent}
\end{equation}
Consequently
\begin{equation}
 J_k-\frac12J_{k-1}
 =b^{k-2}\frac{1-z^{k+1}}{2z(1-z)},
 \qquad
 \mathcal T_k=\frac12b^k\frac{1-z^{k+1}}{1-z}.
 \label{eq:terminal-modal-half-velocity-Laurent}
\end{equation}
Coefficient extraction gives
\[
 t_k(r)=2^{-(k+1)}\sum_{a=|r|}^k\binom ka,
\]
which is \eqref{eq:terminal-modal-half-velocity-tail}.  Convolution on the
$2m$-cycle folds the line coefficient modulo $2m$.  An interior path source
has the two even copies $i$ and $-i$, while the two endpoints have one copy
each.  This proves \eqref{eq:terminal-modal-half-velocity-fold}.  Finally
$q=1/(16m)$ and $qk<3/50$ give $k<48m/50$.
\end{proof}

The same Laurent form retains the derivative cancellation in the proper-face
source ledger.  The following identity is useful because it removes the
apparent geometric sum before any absolute values are taken.

\begin{lemma}[Standard moving-source response; proved here]
\label{lem:terminal-modal-half-velocity-moving-source}
Consider a standard positive-orientation source at transition $n$,
\[
 \widehat F_{n+1}^{U,+}=U_nz^{n-2}(1-z)(1+3z),
 \qquad4\le n\le m-1,
\]
and put $\ell=m-n$.  Its contribution to the full-entry velocity difference
$A_{m+1}-\delta A_m$, followed by $\mathcal T_k$, is exactly
\begin{equation}
 \frac{\delta^\ell U_n}{4}
 b^{k+\ell-1}(1-z^{k+1})(z^{-1}+z^\ell)
 z^{n-2}(1+3z).
 \label{eq:terminal-modal-half-velocity-U-response}
\end{equation}
The reflected source is obtained by $z\mapsto z^{-1}$.  For the accompanying
mass source $M_nz^n$, the corresponding expression is
\begin{equation}
 \frac{\delta^\ell M_n}{4}
 b^{k+\ell-1}\frac{1-z^{k+1}}{1-z}
 (z^{-1}+z^\ell)z^n.
 \label{eq:terminal-modal-half-velocity-M-response}
\end{equation}
The overlapping sources $n<4$ and the final degree-one source at $n=m$ are
not represented by these standard formulas; they remain the exact base and
endpoint groups from
\cref{lem:terminal-modal-boundary-source}.
\end{lemma}

\begin{proof}
A source injected in $A_{n+1}$ contributes
$\delta^{\ell-1}J_\ell F$ to $A_m$ and
$\delta^\ell J_{\ell+1}F$ to $A_{m+1}$.  The elementary identity
\begin{equation}
 J_{\ell+1}-J_\ell
 =\frac12b^{\ell-1}(z^{-1}+z^\ell)
 \label{eq:terminal-modal-entry-velocity-difference}
\end{equation}
therefore gives its contribution to $A_{m+1}-\delta A_m$.
Multiplication by \eqref{eq:terminal-modal-half-velocity-Laurent} cancels
$1-z$ in $\widehat F_{n+1}^{U,+}$ and proves
\eqref{eq:terminal-modal-half-velocity-U-response}.  Without that derivative
factor, the mass source retains the finite quotient in
\eqref{eq:terminal-modal-half-velocity-M-response}.
\end{proof}

\begin{corollary}[Exact finite early certificate; proved here]
\label{cor:terminal-modal-half-velocity-certificate}
For $d$ in \eqref{eq:terminal-modal-static-remainder},
\begin{equation}
 \mathcal L_ku=\mathcal T_kd.
 \label{eq:terminal-modal-half-velocity-static-identity}
\end{equation}
Thus the remaining short-window target
\eqref{eq:terminal-modal-short-early-convolution-open} is exactly the finite
folded inequality
\begin{equation}
 \max_{0\le j\le m}(\mathcal T_kd)(j)\le\frac{q^3}{16},
 \qquad1\le k,\quad qk<\frac3{50},
 \label{eq:terminal-modal-half-velocity-finite-open}
\end{equation}
with the left side given by
\eqref{eq:terminal-modal-half-velocity-fold}.  Equations
\eqref{eq:terminal-modal-half-velocity-U-response} and
\eqref{eq:terminal-modal-half-velocity-M-response} give an exact
source-localized expansion of every standard proper-prefix contribution.
Inequality \eqref{eq:terminal-modal-half-velocity-finite-open} remains
\textbf{OPEN}.
\end{corollary}

\begin{proof}
The deconvolution \eqref{eq:terminal-modal-static-remainder} says $u=Ld$.
Therefore
$\mathcal L_ku=(J_k-J_{k-1}/2)Ld=\mathcal T_kd$.
All remaining statements are substitutions of the preceding two lemmas.
\end{proof}

There is a small but decisive obstruction to a source-free norm shortcut.
At $k=2$ the line coefficients are
\begin{equation}
 t_2(0)=\frac12,qquad t_2(\mathord\pm1)=\frac38,
 \qquad t_2(\mathord\pm2)=\frac18.
 \label{eq:terminal-modal-half-velocity-k2}
\end{equation}
At path target $j=0$, an interior source at $i=1$ contributes both reflected
copies, hence has coefficient $2t_2(1)=3/4$.  Dividing by its degree weight
two gives $3/8>1/4$.  Thus a global claim that every folded coefficient is at
most one quarter per unit $D$-mass is false.  The desired bound must retain
the location of the moving sources and the signed endpoint group; the
positive-mass estimate alone cannot prove it.

The deterministic screen gives
\[
 \max_{j,k:\,qk<3/50}\frac{(\mathcal T_kd)(j)}{q^3}
 =0.059654,0.060267,0.060726,0.061044,0.061268
\]
at $m=64,128,256,512,1024$, respectively.  These values are evidence only.
They are close enough to $1/16$ that discarding reflection or splitting the
final derivative and endpoint atoms by absolute value loses the available
margin.
